\section{Theory of Agent: la «segunda mitad» de la inteligencia de las máquinas}

Wang Hongru (王鸿儒) intenta construir un marco de análisis unificado para los Agent. Los modelos de lenguaje tradicionales optimizan principalmente «qué texto generar después de ver la entrada»; el Agent, en cambio, debe responder «cuándo es necesario actuar, qué capacidad externa debe invocar, cómo mantener un estado a través del tiempo y cómo saber que no sabe».

\sourcefigure{0.93}{lecture08/slides-images/slide-002-00045s.jpg}{Theory of Agent: hacia la segunda etapa de la inteligencia de las máquinas}{00:00:45--00:01:45}

\subsection{Thinking/Reasoning y Doing/Acting}

El informe distingue primero entre el razonamiento interno y la acción externa. Reasoning cambia el token state interno; acting cambia el entorno y obtiene información nueva. Un sistema, aunque pueda producir un chain-of-thought muy largo, no necesariamente sabe cuándo debe detener el razonamiento interno y recurrir a la búsqueda, al código o a la ejecución real.

\sourcefigure{0.93}{lecture08/slides-images/slide-003-00105s.jpg}{Los dos tipos de proceso, thinking y doing, en la inteligencia humana}{00:01:45--00:02:30}

\sourcefigure{0.93}{lecture08/slides-images/slide-004-00150s.jpg}{ReAct organiza reasoning y acting de manera intercalada}{00:02:30--00:03:15}

\begin{importantbox}{La decisión central del Agent}
La acción básica del Agent no es «producir una respuesta», sino elegir entre razonar internamente, invocar una herramienta, observar el entorno, preguntar al usuario y terminar la tarea. La invocación de herramientas es, por lo tanto, un problema de decisión con costo y valor informativo.
\end{importantbox}

\subsection{Frontera del conocimiento y frontera de la capacidad}

El sistema necesita distinguir tres estados: lo que ya se conoce y se puede responder directamente, lo que no se conoce pero se puede obtener mediante una herramienta, y lo que no se conoce y que las herramientas actuales tampoco pueden resolver. Sin un knowledge boundary, el Agent alucina cuando necesita recuperar información; sin un capability boundary, intenta indefinidamente tareas que no puede completar.

\sourcefigure{0.93}{lecture08/slides-images/slide-013-00690s.jpg}{Knowledge Boundary: información conocida, obtenible y no obtenible}{00:11:30--00:13:00}

\sourcefigure{0.93}{lecture08/slides-images/slide-014-00780s.jpg}{La frontera de la capacidad determina cuándo invocar una herramienta y cuándo admitir el fracaso}{00:13:00--00:14:30}

La invocación de herramientas puede escribirse como una comparación de valor:

$$
\mathrm{CallTool}\iff
\mathbb{E}[U\mid \text{tool result}]-C_{\text{tool}}
>\mathbb{E}[U\mid \text{internal answer}],
$$

donde $U$ es la utilidad de la tarea y $C_{\text{tool}}$ incluye la latencia, el costo monetario, la privacidad y el riesgo de ejecución. Lo que el Agent debe aprender es la ganancia neta esperada, y no «usar la herramienta siempre que se pueda».

\subsection{Una perspectiva unificada de contexto largo, RAG, memoria y herramientas}

El contexto largo, la recuperación aumentada, la memoria a largo plazo y la invocación de herramientas amplían la frontera de información que el modelo puede aprovechar en el momento, pero su estructura de costos es distinta. El contexto largo aumenta directamente los tokens; RAG recupera documentos externos; memory conserva el estado relacionado con el usuario o la tarea; las herramientas, en cambio, pueden leer o modificar el entorno.

\sourcefigure{0.93}{lecture08/slides-images/slide-015-00870s.jpg}{La manera en que contexto largo, RAG, memory y herramientas amplían la frontera del conocimiento}{00:14:30--00:15:45}

\sourcefigure{0.93}{lecture08/slides-images/slide-016-00945s.jpg}{La frontera unificada de la capacidad del Agent y la selección de recursos externos}{00:15:45--00:17:45}

\begin{warningbox}{Más context no equivale a mejores decisiones}
Cuanto más largo es el contexto, mayor es el ruido de recuperación y la competencia por la atención; cuantas más herramientas hay, mayor es el costo de elegir mal y el costo de protocolo. El sistema debe aprender resource routing, en lugar de ampliar la entrada sin condiciones.
\end{warningbox}

\subsection{Agentic Post-Training y el objetivo unificado}

Si se quiere que el modelo aprenda cuándo pensar y cuándo actuar, es necesario entrenarlo sobre trayectorias que incluyan herramientas y entorno. Agentic post-training no solo recompensa la respuesta final, sino que también debe abordar el costo de la acción, las restricciones de proceso, la obtención de información y el momento de terminar.

\sourcefigure{0.93}{lecture08/slides-images/slide-019-01425s.jpg}{Agentic post-training: la relación entre SFT, RL y prompt/harness}{00:23:45--00:25:15}

\sourcefigure{0.93}{lecture08/slides-images/slide-020-01515s.jpg}{La Agent policy bajo la perspectiva unificada de post-training}{00:25:15--00:25:30}

\subsection{La formalización del Self-Evolving Agent}

El informe expresa la autoevolución como una actualización de estado a través de episode. Sea $z_t$ el estado interno persistente del Agent, $\tau_t$ la trayectoria de una única interacción y $e_t$ la evaluación; entonces

$$
z_{t+1}=\mathcal{L}(z_t,\tau_t,e_t),
\qquad
\pi_{t+1}(a\mid s)=\pi(a\mid s,z_{t+1}).
$$

$z$ puede incluir los parámetros del modelo, memory, skill, tool policy o harness. Lo esencial de esta definición es que la actualización debe persistir a través de las tareas y poder verificarse en el comportamiento futuro.

\sourcefigure{0.93}{lecture08/slides-images/slide-022-01590s.jpg}{La forma de actualización a través de episode del Self-evolving Agent}{00:26:30--00:27:15}

\sourcefigure{0.93}{lecture08/slides-images/slide-024-01710s.jpg}{La hipótesis de un mundo estático frente al cambio continuo en un mundo abierto}{00:28:30--00:29:15}

\begin{knowledgebox}{El valor de la teoría}
Un marco unificado no necesariamente produce de inmediato una puntuación de benchmark más alta, pero obliga a los investigadores a precisar dónde reside el estado, cuál es el costo de las herramientas, qué escala temporal abarca el aprendizaje, qué variables pueden actualizarse y de dónde provienen las señales de evaluación. Sin estas definiciones, «Agent», «memoria» y «autoevolución» corren el riesgo de quedar reducidos a simples etiquetas de producto.
\end{knowledgebox}

\subsection{Resumen de la sección}

Theory of Agent describe la segunda mitad de la inteligencia de las máquinas como una decisión de recursos y de acción: el modelo debe reconocer la frontera del conocimiento, comparar el valor del razonamiento interno frente al de las herramientas externas, mantener un estado a través del tiempo y, mediante agentic post-training, aprender estrategias de interacción sensibles al costo. Self-evolution es, entonces, la actualización verificable de estos estados persistentes bajo el efecto de la retroalimentación.
